\scanpage{25}
% archive foliation: 24
\uline{Cas $j=i-1$}
\diagram{\begin{tikzcd}[column sep=large,row sep=large]
&\Psi^{i-1}\arrow[r,dash]\arrow[d,"n\operatorname{id}"']&\Phi^i\arrow[r]\arrow[d]&\Psi^i\arrow[d,equal]\\
0\arrow[r,dash]&\Psi^{i-1}\arrow[r]\arrow[d,"\unsure{\theta|0}"']&\Phi(i,n)\arrow[r]\arrow[d,"\alpha=?"]&\Psi^i\arrow[r,dash]\arrow[d,"0"]&0\\
0\arrow[r]&\Psi^{i\deleted{\unsure{+1}}-2}\arrow[r]&\Phi(i-1,m)\arrow[r,dash]&\Psi^{i\deleted{\unsure{+1}}-1}\arrow[r,dash]&0
\end{tikzcd}}
la composée
\diagram{\begin{tikzcd}[column sep=small]
\Phi^i\arrow[r]&\Phi(i,n)\arrow[r,dash]&\Phi(i-1,m)\arrow[r,dash]&\Psi^{i-1}
\end{tikzcd}}
est nulle car $\operatorname{Hom}(\Phi^i,\Psi^{i-1})=0$,
donc l'hom
\deleted{$\Psi^i\longrightarrow\Psi^{\unsure{i-1}}$}
\inserted{$\theta=n\alpha?$}:
$\Psi^{i-1}\longrightarrow\Psi^{i-1}$
composé avec $n\operatorname{id}$ est nul, i.e.
$n\theta=0$, or $\operatorname{Hom}(\Psi^{i-1},\Psi^{i-1})\simeq k$, donc
$\theta=0$, donc on \unsure{reprend} la structure
d'extension, et on termine comme
plus haut pour conclure
$\alpha=0$. Donc on trouve
\[
\operatorname{Hom}(\Phi(i,n),\Phi(j,m))\simeq
\left\{\begin{array}{ll}
0&\text{si }j\neq i,i+1\\
k&\text{si }j=i+1\quad\text{(par la composée d'hom can.)}\\
k&\text{si }j=i\quad\text{(un hom }\deleted{\unsure{respecte}}\text{ ces struct. d'ext.)}
\end{array}\right.
\]
\margin{\diagram{\begin{tikzcd}[column sep=small]
\Phi(i,n)\arrow[r,dash]&\Psi^i\arrow[r]&\Phi(i+1,m)
\end{tikzcd}}}
et est déterminé par ce qu'il induit
sur $\Psi^i$ et $\Psi^{i-1}$, qui doit être de
la forme $t\unsure{m'|n'}\operatorname{id}_{\Psi^i}$ et $t\unsure{n'|m'}\operatorname{id}_{\Psi^{i-1}}$,
si $n'=\dfrac{n}{(m,n)}$, $m'=\dfrac{m}{(m,n)}$ (il y a
une base can. une fois qu'on a
choisi le pgcd $(m,n)$ \unsure{de $m$ et $n$|à un signe}).
